\documentstyle[%


righttag%


,11pt%


,amssymb%


,russian%               remove in Eng.


,izvru%                 Set izven% in Eng.


]{amsart}





%       Math definitions


 


\newcommand{\spc}[1]{{\bold #1}}


\newcommand{\moddd}{\;(\mathrm{mod}\,\mu)\;}


\newcommand{\f}{\varphi}


\newcommand{\e}{\varepsilon}


\newcommand{\tts}[1]{}


\renewcommand{\baselinestretch}{1.06}





\theoremstyle{plain}


\theorembodyfont{\it}


\newtheorem{thmmain}{\hskip\parindent Теорема (основная)}


\renewcommand{\thethmmain}{}





%\hoffset=-15mm%                  For Eng.


 


\setcounter{page}{50}


\god{1999}


\nomer{6~(445)}


\udk{517.988}


\author{\it Ю.В.\,Непомнящих, А.В.\,Поносов}


\title{Локальные операторы в некоторых подпространствах пространства


$L_0$}


\thanks{Работа выполнена при финансовой


поддержке гранта для молодых ученых


(Госкомвуз РФ, 1997) и при частичной финансовой поддержке Российского


фонда фундаментальных исследований, гранты \No\,96-01-01613,


96-15-96195.}





\date{}


%\pagestyle{myheadings}%         Set in Eng.


\pagestyle{plain} %              Remove in Eng.


\begin{document}


%\thispagestyle{plain}%          Set in Eng.


\maketitle


\label{begin}





\section{Введение}





Теория локальных операторов (ЛО) --- одно из перспективных направлений


современного нелинейного анализа, которое бурно развивается на


протяжении второй половины двадцатого века и находит многочисленные


приложения в механике, физике, теории управления, других науках.


Понятие ЛО возникло путем абстрагирования от известного свойства


оператора Немыцкого (су\-пер\-по\-зи\-ции) и заключается в том, что значение


функ\-ции--об\-ра\-за на некотором множестве зависит только от значений


функ\-ции--про\-об\-ра\-за на том же множестве (точ\-ное определение


см. ни\-же). Впервые определение ЛО в его современном понимании


дано в \cite{1} (там они названы ``ло\-каль\-но определенными


опе\-ра\-то\-ра\-ми''), хотя близкие к этому понятия встречались и ранее


(\cite{2}, \cite{3}).





Фундаментальным является то обстоятельство, что непрерывный по мере ЛО


в ``клас\-си\-чес\-ких'' пространствах измеримых функций представим в


виде оператора Немыцкого с порождающей функцией, удовлетворяющей


условиям Каратеодори. Этот факт в некоторых частных случаях доказан


в \cite{3}, а в общем случае --- в \cite{4}--\cite{6}.


Более


конструктивное доказательство теоремы о представлении (в менее общей


по сравнению с \cite{4}--\cite{6} ситуации) было


предложено в \cite{7},\cite{8}. Таким образом,


непрерывный ЛО есть оператор Немыцкого с порождающей функцией,


удовлетворяющей условиям Каратеодори --- объект, к настоящему времени


изученный достаточно хорошо с разных сторон (напр.,


\cite{1}--\cite{13}, в обзорной монографии \cite{8} имеется обширная


библиография).





В работах \cite{14}, \cite{15} было замечено, что свойство локальности


присуще типичным операторам, возникающим в теории стохастических


дифференциальных и интегральных уравнений. В стохастике характерна ситуация,


когда ЛО исследуется в пространствах функций, измеримых по отношению к


различным


$\sigma$-алгебрам на области определения этих функций. Разумеется,


здесь речь идет о ситуации, когда функции из области задания ЛО


измеримы относительно более узкой $\sigma$-ал\-геб\-ры, чем функции из


множества значений оператора. Оказывается, что в такой


``не\-клас\-си\-чес\-кой'' постановке непрерывные по мере ЛО не


сводятся к операторам Немыцкого. Характерный пример --- стохастические


интегральные операторы (в том числе и не\-ли\-ней\-ные).





Другой мотивацией для изучения ЛО являются совсем недавние исследования


М.Е.\,Драх\-лина, Е.\,Степанова и второго соавтора (статья


``Memory properties of operators'' представлена к рассмотрению для


публикации в зарубежном издании), показывающие,


что задача классификации так называемых


``ато\-ми\-чес\-ких''


операторов сводится по сути к изучению ЛО, которые заданы на


подпространствах пространства $L_0$, состоящих из функций, измеримых


по отношению к $\sigma$-подал\-геб\-ре заданной $\sigma$-ал\-геб\-ры. Классу


атомических операторов принадлежат, кроме ЛО, многие операторы с


``па\-мя\-тью'', в частности, различные композиции подстановки (опе\-ра\-то\-ра


внутренней суперпозиции \cite{20}, \cite{33}) и ЛО $h$. Изучение таких


операторов в свою очередь находит непосредственное применение для


исследования функ\-ци\-о\-наль\-но-диф\-фе\-рен\-ци\-аль\-ных уравнений


(см., напр., \cite{33}, \cite{34}).





В данной работе, как нам представляется, впервые


изучаются ЛО в описанной выше ``не\-клас\-си\-чес\-кой''


постановке. Более конкретно,


выявлено условие на $\sigma$-под\-ал\-геб\-ру,


необходимое и достаточное для представления


непрерывного по мере ЛО в виде оператора Немыцкого с порождающей


функцией, удовлетворяющей условиям Каратеодори.


Подчеркнем, что в работе мы изучаем только непрерывные ЛО.


Без условия непрерывности фактически любой ЛО есть


оператор Немыцкого (\cite{22}, \cite{7}, \cite{8}, \cite{23}), однако


порождающая функция может обладать весьма патологическими свойствами


(так называемые ``уродцы''). Заметим лишь, что в указанных работах факт


существования ``уродцев'' доказан в предположении


кон\-ти\-ну\-ум-ги\-по\-те\-зы (CH), хотя на самом деле (\cite{24}, c.\,145)


достаточно


аксиомы Мартина (MA) (напомним, что (CH)$\Rightarrow$(MA)


и (MA)\&$\neg$(CH) не противоречит аксиомам формальной теории


множеств). В предположении MA существует такая неизмеримая (по Лебегу)


функция $f:[0,1]^2\to{\spc R}$, что для любой измеримой функции


$\f:[0,1]\to[0,1]$ при почти всех $t\in [0,1]$ имеет место равенство


$f(\f(t),t)=f(t,\f(t))=0$.





\section{Локальные операторы на склейках функций}





Пусть $(T,\Sigma)$ --- измеримое пространство, $\Sigma_0$ ---


$\sigma$-пол\-ный идеал алгебры $\Sigma$ \cite{25}; $X$, $Y$ ---


произвольные непустые множества.\; ${\cal F}(E,X):=X^E$


$(E\subset T)$ --- множество


всех отображений из $E$ в $X$. На ${\cal F}(E,X)$


рассмотрим естественное отношение эквивалентности $\cal R$:


$[\f\sim\psi] \ \Leftrightarrow \


[\{t\mid\f(t)=\psi(t)\}\in\Sigma_0]$.\; $F(E,X):={\cal F}(E,X)/{\cal R}$


--- соответствующее фак\-тор-про\-ст\-ран\-ст\-во


классов эквивалентных функций из $E$ в $X$. Для краткости положим


${\cal F}(X):={\cal F}(T,X)$, $F(X):= F(T,X)$.





Определим оператор проектирования ${\cal P}_{E,X}:{\cal F}(X)\to {\cal


F}(E,X)$ равенством ${\cal P}_{E,X}\f:=\f|_E$ (сужение функции $\f$ на


множество $E$). Очевидно, корректно определен соответствующий оператор


проектирования на фак\-тор-про\-ст\-ран\-ст\-вах $P_{E,X}:F(X)\to


F(E,X)$. Далее часто для краткости вместо ${\cal P}_{E,X}$, $P_{E,X}$


будем писать ${\cal P}_{E}$, $P_{E}$. Из контекста всегда ясно, о


функциях с какой областью значений идет речь.





Пусть $R=\{T_n\}_{n\in I}$ --- счетное разбиение пространства $T$.


Здесь и далее разбиение


называем счетным, если множество индексов $I$ непусто и либо конечно,


либо счетно.





\begin{defn}


{\it Склейкой} множества $M\subset F(X)$,


соответствующей разбиению $R$ (обозначение $S_R(M)$), назовем множество


$$


S_R(M):=\{\f\in F(X)\mid(\exists \f_n\in M) \


P_{T_n}\f_n=P_{T_n}\f, \ n\in I\}.


$$


\end{defn}





\begin{lem}


Если $\f,\psi\in F(X)$ и $P_{T_n}\f=P_{T_n}\psi$,


$n\in I$, то $\f=\psi$.


\end{lem}





Доказательство вытекает непосредственно из $\sigma$-пол\-но\-ты идеала


$\Sigma_0$.





\begin{lem}


Пусть $M,\,M_1,\,M_2\in F(X)$. Имеют место


свойства


$$


[M_1\subset M_2]\Rightarrow[S_R(M_1)\subset S_R(M_2)], \quad


M\subset S_R(M), \quad S_R(S_R(M))=S_R(M).


$$


\end{lem}





\begin{pf} Первые два свойства очевидны. Для доказательства


последнего равенства зафиксируем произвольное $\f\in S_R(S_R(M))$


и найдем согласно определению $S_R(S_R(M))$ такие $\f_n\in


S_R(M)$, что $P_{T_n}\f_n=P_{T_n}\f$, $n\in I$.


Далее, найдем согласно определению $S_R(M)$


такие $\psi_n\in M$, что $P_{T_n}\psi_n=P_{T_n}\f_n$, $n\in I$. Тогда


$P_{T_n}\psi_n=P_{T_n}\f$, $n\in I$, и, таким образом, $\f\in S_R(M)$.


Итак, доказано $S_R(S_R(M))\subset S_R(M)$.


Обратное включение вытекает из первых двух свойств этой леммы.


\end{pf}





\begin{defn}


Оператор $h:M\to F(Y)$ ($M\subset F(X)$)


называется {\it локальным оператором} \cite{1} (сокращенно --- ЛО),


если


$$


[\f,\psi\in M, \ E\in\Sigma, \ P_{E}\f=P_{E}\psi


]\Longrightarrow[P_{E}h\f=P_{E}h\psi].


$$


\end{defn}





\begin{thm}


Пусть $M\subset F(X)$, и $h:M\to F(Y)$ --- ЛО.


Тогда существует единственный ЛО $\wt{h}:S_R(M)\to F(Y)$,


для которого $\wt{h}|_M=h$. При этом $\wt{h}(S_R(M))=S_R(h(M))$.


\end{thm}





\begin{pf} 1) Сначала докажем единственность ЛО


$\wt{h}:S_R(M)\to F(Y)$, совпадающего с $h$ на $M$.


Пусть $h_k$, $k=1,2$, --- два таких продолжения. Для любого


$\f\in S_R(M)$ и для любого $n\in I$ найдем такое


$\f_n\in M$, что $P_{T_n}\f=P_{T_n}\f_n$. Из равенства $h_k|_M=h$


и локальности операторов $h_k$ следует


$$


P_{T_n}h_1\f=P_{T_n}h_1\f_n=P_{T_n}h\f_n=


P_{T_n}h_2\f_n=P_{T_n}h_2\f, \quad n\in I.


$$





В силу леммы 1\; $h_1\f=h_2\f$, и ввиду произвольности


$\f\in S_R(M)$ имеем $h_1=h_2$.





2) Для произвольного $\f\in S_R(M)$


найдем такое множество $\{\f_n\}_{n\in I}\subset M$, что


$P_{T_n}\f=P_{T_n}\f_n$, $n\in I$, причем если $\f\in M$, то считаем


$\f_n=\f$. Определим оператор $\wt{h}$ равенством


\begin{equation}


\label{1}


(\wt{h}\f)(t)=(h\f_n)(t)\;\ text{почти всюду (п.\,в.) на}\; \ T_n,


\quad n\in I.


\end{equation}


Очевидно, $\wt{h}|_M=h$,


$\wt{h}(S_R(M))\subset S_R(h(M))$.  Итак, построено продолжение


$\wt{h}$ оператора $h$ на $S_R(M)$.





3) Покажем, что $\wt{h}$ --- ЛО. Пусть $\f,\psi\in S_R (M)$,


$E\subset T$, $P_{E}\f=P_{E}\psi$. Определим $\f_n,\psi_n\in M$, $n\in I$,


так, что $P_{T_n}\f=P_{T_n}\f_n$, $P_{T_n}\psi=P_{T_n}\psi_n$,


$n\in I$, выполнено равенство (\ref{1}) и равенство


\begin{equation}


\label{1111}


(\wt{h}\psi)(t)=(h\psi_n)(t)\;\ \text{п.\,в. на}\; \ T_n, \quad n\in I.


\end{equation}


Очевидно, $P_{E\cap T_n}\f_n=P_{E\cap T_n}\f=P_{E\cap T_n}\psi=


P_{E\cap T_n}\psi_n$, и согласно (\ref{1}), (\ref{1111}) и локальности $h$


$$


P_{E\cap T_n}\wt{h}\f=P_{E\cap T_n}h\f_n=


P_{E\cap T_n}h\psi_n=P_{E\cap T_n}\wt{h}\psi, \quad n\in I.


$$


В силу леммы 1 $P_E\wt{h}\f=P_E\wt{h}\psi$.





Осталось показать $S_R(h(M))\subset\wt{h}(S_R(M))$.


Фиксируем произвольное $\psi\in S_R(h(M))$ и выбираем $\psi_n\in


h(M)$ такие, что $P_{T_n}\psi_n=P_{T_n}\psi$, $n\in I$. Найдем


$\f_n\in M$, для которых $\psi_n=h\f_n$, $n\in I$. Определим


$\f\in S_R(M)$, положив $\f(t)=\f_n(t)$, $t\in T_n$ $(n\in I)$.


В силу локальности $\wt{h}$ имеем


$$


P_{T_n}\wt{h}\f=P_{T_n}\wt{h}\f_n=


P_{T_n}h\f_n=P_{T_n}\psi_n=P_{T_n}\psi, \quad n\in I.


$$


Согласно лемме 1\; $\psi =\wt{h}\f\in\wt{h}(S_R(M))$.


\end{pf}





\begin{note}


Роль инвариантных относительно


склейки множеств для понимания природы ЛО


впервые замечена, по-ви\-ди\-мо\-му, в работах \cite{1}, \cite{26}


(см. также \cite{22}, \cite{23}). Эти


понятия тесно связаны с определенным и изученным в \cite{8} понятием


``thick set''.


\end{note}





\begin{note}


Обобщая результаты этого параграфа, можно дать


определение склей\-ки $S_{\cal A}(M)$, соответствующей некоторому


семейству разбиений $\cal A$ (например, множеству всех счетных


измеримых разбиений). При естественных ограничениях на


$(S,\Sigma)$ и $\cal A$ в этом случае справедлив аналог теоремы 1.


Это свойство ЛО (в менее абстрактной постановке и несколько в других


терминах) указывалось и использовалось в работах \cite{26}, \cite{8}.


Мы не приводим соответствующего


обобщения, поскольку степень общности в теореме 1 достаточна для


доказательства основных результатов статьи.


\end{note}





\section{Некоторые свойства $\sigma$-алгебр и пространства


$L_0(\Sigma, X)$}





Пусть $(T,\Sigma,\mu)$ --- пространство с полной конечной мерой,


$(X,d)$ --- сепарабельное метрическое пространство, ${\cal


B}(X)$ --- $\sigma$-ал\-геб\-ра борелевских подмножеств $X$. Функцию


$\f:T\to X$ назовем $\Sigma$-из\-ме\-ри\-мой, если $(\forall U\in{\cal


B}(X))$\; $\f^{-1}(U)\in\Sigma$. Всюду $L_0(\Sigma,X)$ --- множество классов


эк\-ви\-ва\-лент\-ных из\-ме\-ри\-мых функций из $T$ в $X$, наделенное


топологией сходимости по мере.





Пусть $E\subset T$ --- фиксированное (вообще говоря, неизмеримое)


множество. Тогда $(E,\Sigma\cap E,\mu_E)$ ---


пространство с полной конечной мерой. Здесь и далее $\Sigma\cap


E:=\{A\cap E\mid A\in\Sigma\}$ \cite{27}, $\mu_E$ ---


сужение внешней меры $\mu^\ast$ на $\Sigma\cap E$. $L_0(\Sigma\cap


E,X)$ --- пространство $\Sigma\cap E$-из\-ме\-ри\-мых функций из $E$ в


$X$.





Для произвольного $A\in\Sigma\cap E$ положим ${\cal E}(A):=


\{B\in\Sigma\mid B\cap E=A\}$, $\ov{\cal E}(A):=\{B\in\Sigma\mid


B\supset A\}$. ${\cal E}(A)$, $\ov{\cal E}(A)$ состоят из


элементов булевой фак\-тор-ал\-геб\-ры $\Sigma/\Sigma_0$ ($B_1=B_2


\moddd$, если $\mu(B_1\bigtriangleup B_2)=0$).


На $\Sigma/\Sigma_0$ рассматриваем естественное отношение


частичного порядка (т.\,е. по включению $\moddd$).


Часто в дальнейшем,


допуская некоторую вольность, будем с элементами $\Sigma/\Sigma_0$


обращаться как с множествами, считая, что равенства и включения


выполняются $\moddd$. Это не приведет к недоразумениям.





\begin{lem}


${\cal E}(A)$, $\ov{\cal E}(A)$ имеют


наименьшие элементы, равные между собой.


\end{lem}





\begin{pf} Пусть $a=\inf\{\mu B\mid B\in{\cal E}(A)\}$, и $B_n\in{\cal


E}(A)$ таковы, что $\mu B_n-a<1/n$, $n=1,2,\dotsc$ Тогда $A^\ast:=


\mathop{\cap}\limits_{n=1}^\infty B_n\in{\cal E}(A)$ и $\mu A^\ast=a$.





Предположим, что $A^\ast$ не является наименьшим элементом ${\cal


E}(A)$. Тогда существует $A_1\in{\cal E}(A)$


такое, что $\mu(A^\ast\setminus A_1)>0$. В этом случае $A_2:= A^\ast\cap


A_1\in{\cal E}(A)$, $\mu A_2<a$, что противоречит определению числа


$a$.





Аналогично доказывается, что существует наименьший элемент


$\ov{\cal E}(A)$. Обозначим его символом $\ov{A}$.


Из очевидного включения ${\cal E}(A)\subset\ov{\cal E}(A)$


следует $\ov{A}\subset A^\ast$, поэтому


$A\subset\ov{A}\cap E\subset A^\ast\cap E=A$. Таким образом,


$\ov{A}\cap E=A$, откуда следует $\ov A\in{\cal


E}(A)$. Согласно определению $A^\ast$ имеет место


$A^\ast\subset\ov{A}$. Равенство $A^\ast=\ov{A}\moddd$


доказано.


\end{pf}





Положим $A^\ast:=\inf{\cal E}(A)=\inf\ov{\cal E}(A)$.





\begin{lem}


Отображение $\gamma$, определяемое


соотношением $A\mapsto A^\ast$, есть изоморфизм булевой алгебры


$(\Sigma\cap E)/\Sigma_0$ на алгебру $(\Sigma\cap E^\ast)/\Sigma_0$.


\end{lem}





\begin{pf} Последовательно докажем ряд свойств.





1) $[A,B\in\Sigma\cap E,\ A\subset B]\Rightarrow[\gamma(A)


\subset\gamma(B)]$ (в частности, $\gamma(A)\in(\Sigma\cap


E^\ast)/\Sigma_0$).





Очевидно, $\gamma(B)\in\ov{\cal E}(A)$, и поскольку (согласно


лемме 3) $\gamma(A)$ --- наименьший элемент $\ov{\cal E}(A)$,


то $\gamma(A)\subset\gamma(B)$.





2) $[A_n\in\Sigma\cap E,\ n=1,2,\dotsc]\Rightarrow[


\gamma(\mathop{\cup}\limits_{n=1}^\infty


A_n)=\mathop{\cup}\limits_{n=1}^\infty\gamma(A_n)]$.





Из свойства 1) следует


$\mathop{\cup}\limits_{n=1}^\infty\gamma(A_n)\subset\gamma(\mathop


{\cup}\limits_{n=1}^\infty A_n)$. Далее, очевидно,


$\mathop{\cup}\limits_{n=1}^\infty\gamma(A_n)\in\ov{\cal


E}(\mathop{\cup}\limits_{n=1}^\infty A_n)$, и поскольку


$\gamma\left(\mathop{\cup}\limits_{n=1}^\infty A_n\right)$ есть


наименьший элемент $\ov{\cal


E}\left(\mathop{\cup}\limits_{n=1}^\infty(A_n)\right)$, то


$\gamma\left(\mathop{\cup}\limits_{n=1}^\infty


A_n\right)\subset\mathop{\cup}\limits_{n=1}^\infty\gamma(A_n)$.





3) $\forall B\in\Sigma\cap\gamma(E)$ $\gamma(B\cap E)=B$ (в частности,


$\gamma$ сюръективно).





Пусть $B\in\Sigma\cap\gamma(E)$. Тогда $A:= B\cap E\in\Sigma\cap E$.


Согласно лемме 3 и определению семейств ${\cal E}(A)$,


$\ov{\cal E}(A)$


имеем $\gamma(A)\subset B$, $A=\gamma(A)\cap E=B\cap E$. Из последнего


равенства и свойства~1) следует $D:=


B\setminus\gamma(A)\subset\gamma(E)\setminus E$.


Отсюда, из определения $\gamma(E)$ как наибольшего элемента


$\ov{\cal E}(E)$ и включения $D\in\Sigma$ вытекает


$D\in\Sigma_0$. Таким образом, $\gamma(A)=B\moddd$.





4) $[A,B\in\Sigma\cap E,\ A\cap B\in\Sigma_0]\Rightarrow[


\gamma(A)\cap\gamma(B)\in\Sigma_0]$.





Для любого $D\in\Sigma\cap E$ имеет место $\gamma(D)\cap E=D$


(это следует из включения $\gamma(D)\in{\cal E}(D)$). Из


свойства~3) и этого равенства вытекает


$$


\gamma(A)\cap\gamma(B)=\gamma(\gamma(A)\cap\gamma(B)\cap E)=


\gamma(A\cap B)=\gamma(\emptyset)=\emptyset\,\moddd.


$$





5) $[A,B\in\Sigma\cap E]\Rightarrow


[\gamma(E\setminus B)=\gamma(E)\setminus\gamma(B)]$.





Из 2), 4) следует $\gamma(B)\cup\gamma(E\setminus B)=\gamma(E)$,


$\gamma(B)\cap\gamma(E\setminus B)=\emptyset$. Отсюда, очевидно,


вытекает требуемое равенство.





6) $\gamma$ взаимно однозначно отображает $(\Sigma\cap E)/\Sigma_0$ на


$(\Sigma\cap\gamma(E))/\Sigma_0$. При этом


$\gamma^{-1}$ определяется равенством $\gamma^{-1}(B)=B\cap E$.





Если $A\ne B$, то


согласно 3), 5)\; $\gamma(A)\bigtriangleup\gamma(B)=\gamma(A\bigtriangleup


B)\notin\Sigma_0$, т.\,е.


$\gamma(A)\ne\gamma(B)$. Таким образом, отображение $\gamma$


инъективно. Сюръективность $\gamma$ и равенство $\gamma^{-1}(B)=B\cap


E$ вытекают из утверждения п.\,3).


\end{pf}





\begin{thm}


Для любого фиксированного $\wt{x}\in X$


существует единственный оператор


$r_E:L_0(\Sigma\cap E,X)\to L_0(\Sigma ,X)$, для


которого


\begin{equation}


\label{2}


P_Er_E\f=\f, \quad P_{T\setminus\gamma(E)}r_E\f\equiv\wt{x}.


\end{equation}


\end{thm}





\begin{pf} 1) Пусть ${\cal P}_0(\Sigma\cap E,X)$, ${\cal P}_0(\Sigma


,X)$ --- пространства счетнозначных $\Sigma\cap E$-из\-ме\-ри\-мых


функций из $E$ в $X$ и соответственно $\Sigma$-из\-ме\-ри\-мых


функций из $T$ в $X$.





Зафиксируем произвольное $\f\in{\cal P}_0(\Sigma\cap E,X)$


вида $\f(t)=x_k$, $t\in E_k$, где $R:=


\{E_{k}\}_{k\in I}\subset\Sigma$ --- счетное разбиение $E$.


Определим функцию $r_E\f$ равенством


\begin{equation}


\label{4}


(r_E\f)(t)=x_k,\;\ t\in\gamma(E_{k}), \quad


(r_E\f)(t)=\wt{x},\;\ t\in T\setminus\gamma(E)


\end{equation}


($\wt{x}\in X$ произвольно и фиксировано).


Определение корректно, т.\,к. согласно лемме 4


$\{\gamma(E_{k})\}_{k\in I}$ есть счетное


измеримое разбиение $\gamma(E)$.





Тем самым мы определили оператор $r_E:{\cal P}_0(\Sigma\cap


E,X)\to{\cal P}_0(\Sigma,X)$. Для $\f\in{\cal P}_0(\Sigma\cap E,X)$


равенства (\ref{2}) очевидным образом вытекают из построения.





2) Зафиксируем теперь произвольное $\f\in L_0(\Sigma\cap E,X)$ и найдем


такие $\f_n\in{\cal P}_0(\Sigma\cap E,X)$, $n=1,2,\dotsc$, что


$\f_n\to \f$ п.\,в. Для произвольного $\sigma >0$ положим


$G_{nm}=\gamma(A_{nm})$, где $A_{nm}=\{t\in


E\mid d(\f_n(t),\f_m(t))>\sigma\}$. Согласно (\ref{4})\; $G_{nm}=\{t\in


T\mid d((r_E\f_n)(t),(r_E\f_m)(t))>\sigma\}$. Поскольку


$\f_n@>{\mu_E}>>\f$, то в силу леммы 4\; $\mu


G_{nm}=\mu_E A_{nm}\to 0$, $n,m\to\infty$. Таким образом,


последовательность $\{r_E\f_n\}$ фундаментальна по мере, и согласно


секвенциальной полноте пространства $L_0(\Sigma, X)$ она


сходится к некоторому элементу пространства $L_0(\Sigma, X)$, который


обозначим через $r_E\f$. Извлекая из $\{r_E\f_n\}$ сходящуюся почти


всюду подпоследовательность $\{r_E\f_{n_k}\}$, получаем из свойств


$P_Er_E\f_{n_k}=\f_{n_k}$,


$P_{T\setminus\gamma(E)}r_E\f_{nk}\equiv\wt{x}$ и сходимости


$\f_{n_k}@>{\text{п.в.}}>>\f$ справедливость


равенств (3).





Таким образом, оператор $r_E$ на $L_0(\Sigma\cap E,X)$, удовлетворяющий


свойству (\ref{2}), построен.





3) Докажем единственность оператора $r_E$, удовлетворяющего (\ref{2}). Пусть


$\f\in L_0(\Sigma\cap E,X)$, $\psi_i\in L_0(\Sigma, X)$, $i=1,2$ таковы,


что


$$


P_E\psi_i=\f, \quad P_{T\setminus\gamma(E)}\psi_i\equiv\wt{x},


\quad i=1,2.


$$


Очевидно, $B:=\{t\in T\mid\psi_1(t)\ne\psi_2(t)\}=B_1\cup N$, где $\mu


N=0$, а $B_1\subset\gamma(E)\setminus E$. По лемме 4\;


$\mu^\ast(\gamma(E)\setminus E){=}0$, поэтому $\mu^\ast B_1{=}0$.


В силу полноты меры $\mu B_1{=}0$. Таким образом, $\psi_1{=}\psi_2$.


\end{pf}





\section{$\Omega$-условие и его свойства}





Пусть $(T,\Sigma,\mu)$ --- пространство с полной конечной мерой,


$(X,d)$, $(Y,\rho)$ --- полные сепарабельные метрические пространства,


$\Sigma_1\subset\Sigma$ --- $\sigma$-под\-ал\-геб\-ра алгебры $\Sigma$,


полная относительно $\mu$. Последнее означает


$\Sigma_1\supset\Sigma_0:=\mu^{-1}(\{0\})$.





\begin{defn}


Будем говорить, что $\sigma$-по\-дал\-геб\-ра


$\Sigma_1$\; $\sigma$-алгебры $\Sigma$ удовлетворяет


$\Omega$-{\it ус\-ло\-вию} (со\-кра\-щен\-но $\Sigma_1\in{\Omega}


(\Sigma)$), если


существует счетное измеримое разбиение $R=\{T_n\}_{n\in I}$


пространства $T$ со свойством


\begin{equation}


\label{5}


\Sigma_1\cap T_n=\Sigma\cap T_n, \quad n\in I.


\end{equation}


\end{defn}





Для $\Sigma_1\in{\Omega}(\Sigma)$ через $\Re(\Sigma_1)$ будем обозначать


семейство счетных измеримых разбиений $\moddd$ $R=\{T_n\}_{n\in I}$,


удовлетворяющих (\ref{5}).





\begin{note}


В \cite{20}


определяется условие $\omega(h_1,h_2,\dotsc,h_n)$ на функцию


$g:{\spc R}^m\to{\spc R}^m$ (случай $n=\infty$ не исключается), при


выполнении которого $\sigma$-по\-дал\-геб\-ра исходной


$\sigma$-ал\-геб\-ры в ${\spc R}^n$ как раз удовлетворяет


$\Omega$-ус\-ло\-вию в смысле определения~3. В \cite{20} условие


$\omega$ играет ключевую роль для нахождения аналитического представления


оператора, сопряженного к линейному оператору внутренней


суперпозиции. Примечательно, что близкое


$\Omega$-ус\-ло\-вие играет не менее важную роль в вопросе о


представлении нелинейного ЛО в виде оператора Немыцкого (см.


\S\,5 ниже).


\end{note}





Зафиксируем произвольную $\sigma$-под\-ал\-геб\-ру $\Sigma_1$


алгебры $\Sigma$ и рассмотрим ее строение. Всюду далее считаем,


что элементы $\Sigma/\Sigma_0$ упорядочены естественным образом,


т.\,е. по включению $\moddd$.





\begin{lem}


Если семейство ${\cal E}\subset\Sigma/\Sigma_0$ замкнуто


относительно счетных объединений, то оно имеет наибольший элемент.


\end{lem}





\begin{pf} Пусть $a=\sup\limits_{E\in{\cal E}}\mu E$.


Найдем $E_n\in{\cal E}$, для которых $\mu E_n>a-1/n$, $n=1,2,\dotsc$, и


положим $E=\mathop{\cup}\limits_{n=1}^\infty E_n$. По условию


$E\in{\cal E}$, и согласно построению $\mu E=a$. Покажем, что $E$ ---


наибольший элемент $\cal E$.





Предположим противное, т.\,е. что $(\exists D\in{\cal E})$


$\mu(D\setminus E)>0$. Тогда $D\cup E\in\Sigma$, и $\mu(D\cup E)>a$,


что противоречит определению числа~$a$.


\end{pf}





\begin{lem}


Семейство


$\Sigma_c:=\{E\in\Sigma_1\mid \Sigma_1\cap E=\Sigma\cap E\}/\Sigma_0$


имеет наибольший элемент.


\end{lem}





{\bf Доказательство.} В силу леммы 5 достаточно показать, что $\Sigma_c$


замкнуто относительно счетных объединений. Пусть $E_n\in\Sigma_c$,


$n=1,2,\dotsc$, $E=\mathop{\cup}\limits_{n=1}^\infty E_n$, и


$A\in\Sigma\cap E$ произвольно. Поскольку $E_n\in\Sigma_c$, то $A_n:=


A\cap E_n\in\Sigma_1\cap E_n\subset\Sigma_1$. Тогда


$A=\mathop{\cup}\limits_{n=1}^\infty A_n\in\Sigma_1$. Так как


вдобавок


$A\subset E$, то  $A\in\Sigma_1\cap E$.


Включение $E\in\Sigma_c$ доказано.


%$\quad\Box$





Наибольший элемент семейства $\Sigma_c$, который существует


в силу леммы 4, обозначим символом $T_c$.





\begin{lem}


Семейство


$\Sigma_d:=\{E\in\Sigma\cap(T\setminus T_c)\mid$ существует такое счетное


измеримое разбиение $\{E_n\}_{n\in I}$ множества $E$, что


$\Sigma_1\cap E_n=\Sigma\cap E_n,\ n\in I\}/\Sigma_0$


имеет наибольший элемент.


\end{lem}





\begin{pf} Как и в доказательстве леммы 6, достаточно доказать


замкнутость $\Sigma_d$ относительно счетных объединений.





1) Проверим свойство


\begin{equation}


\label{6}


[E\in\Sigma_d,\ A\in\Sigma\cap E]\,\Rightarrow\,[\,


A\in\Sigma_d].


\end{equation}





Действительно, если $E\in\Sigma_d$ и $\{E_n\}_{n\in I}$ --- счетное


измеримое разбиение $E$, для которого $\Sigma_1\cap E_n=\Sigma\cap


E_n$, $n\in I$, то для элементов $A_n:= E_n\cap A$ счетного


измеримого разбиения $\{A_n\}_{n\in I}$ множества $A$ справедливо


равенство $\Sigma_1\cap A_n=\Sigma\cap A_n$.


Включение $A\in\Sigma_d$ доказано.





2) Пусть теперь $E_n\in\Sigma_d$, $E=\mathop{\cup}\limits_{n=1}^\infty


E_n$, и счетные измеримые разбиения $\{T_{nk}\}_{k\in I_n}$ множеств


$E_n$, $n=1,2,\dotsc$, выбраны так, что $\Sigma_1\cap T_{nk}=\Sigma\cap


T_{nk}$ ($\forall k,n$).





Положим $A_1=E_1$,


$A_n=E_n\setminus\mathop{\mathop{\cup}\limits}\limits_{k=1}^{n-1}


E_k$, $k=2,3,\dotsc$ Очевидно, $E=\mathop{\cup}\limits_{n=1}^\infty A_n$


и множества $A_n$ попарно дизъюнктны. В силу (\ref{6})\; $A_n\in\Sigma_d$.





Пусть измеримые разбиения


$\{A_{nk}\}_{k\in I_n}$ множеств $A_n$, $n=1,2,\dotsc$,


выбраны так, что


\begin{equation}


\label{7}


\Sigma_1\cap A_{nk}=\Sigma\cap A_{nk}, \quad k\in I_n,\;\ n=1,2,


\dotsc


\end{equation}


Очевидно, $\{A_{nk}\}_{k\in I_n,\,n\in I}$ ---


счетное измеримое разбиение $E$, и в силу (\ref{7}) имеем $A\in\Sigma_d$.


\end{pf}





Наибольший элемент семейства $\Sigma_d$ обозначим символом $T_d$.


Положим $T_s=T\setminus(T_c\cup T_d)$.





\begin{thm}


Семейство $T$ однозначно $\moddd$ представимо


в виде дизъюнктного объединения трех множеств$:$ $T_c,\,T_d,\,T_s\in\Sigma$,


которые обладают свойствами


\begin{itemize}


\item[1)] $T_c\in\Sigma_c;$ для любого $D\in\Sigma_1$, удовлетворяющего


условию $\Sigma_1\cap D=\Sigma\cap D$, имеет место включение $D\subset


T_c\moddd;$


\item[2)] $T_d\in\Sigma_d$ и $(\forall D\in\Sigma_d)$\; $D\subset T_d\moddd$,


более того, $(\forall D\in(\Sigma_1\cap T_d)\setminus\Sigma_0)$\;


$\Sigma_1\cap D\ne\Sigma\cap D\moddd$, в частности, элементы любого


счетного измеримого разбиения $\{T_n\}_{n\in I}$ множества $T$,


удовлетворяющего условию $\Sigma_1\cap T_n=\Sigma\cap T_n$, $n\in I$,


не принадлежат $\Sigma_1;$


\item[3)] $(\forall D\in(\Sigma\cap T_s)\setminus\Sigma_0)$\;


$\Sigma_1\cap D\ne\Sigma\cap D$.


\end{itemize}


\end{thm}





\begin{pf} Все утверждения теоремы очевидным образом вытекают


из лемм 6, 7. Заметим лишь, что если предположить для некоторого


$D\in\Sigma_1\cap T_d$ справедливость $\Sigma_1\cap D=\Sigma\cap D$,


то получим


$T_c\cup D\in\Sigma_c$ в противоречие с определением $T_c$, а если бы


нашлось $D\in\Sigma\cap T_s$, для которого $\Sigma_1\cap D=\Sigma\cap D$,


то $D\in\Sigma_c\cup\Sigma_d$ в противоречие


с определениями $T_c$ и $T_d$.


\end{pf}





\begin{cor}


1) $[T_c=T\moddd]\Leftrightarrow[\Sigma=\Sigma_1]$;





2) $[T_s=\emptyset]\Leftrightarrow[\Sigma_1\in{\Omega}(\Sigma)]$;





3) $[T_s=T\moddd]\Leftrightarrow[(\forall E\in\Sigma)\;


\Sigma_1\cap E\not\in{\Omega}(\Sigma\cap E)]$;





4) $T_s$ не содержит атомов $\sigma$-алгебры $\Sigma$.


\end{cor}





Из вторых утверждений теоремы 3 и следствия 1 получаем


следующий критерий.





\begin{cor}


$\sigma$-подалгебра $\Sigma_1$


удовлетворяет $\Omega$-ус\-ло\-вию тогда и только тогда, когда


$$


(\forall E\in\Sigma)\;\ (\exists D\in\Sigma\cap E):\;\


\Sigma_1\cap D=\Sigma\cap D.


$$


\end{cor}





\begin{note}


По сути большинство результатов этого параграфа


следует из более общих результатов теории меры (\cite{25},


\cite{28}, \cite{29})


и допускает разнообразную интерпретацию.


В частности, $\Omega$-ус\-ло\-вие означает, что в $T$ нет


$\Sigma_1$-од\-но\-род\-ных ненасыщенных множеств (\cite{29},


гл.\,2) --- об этом говорит следствие 2.


Однако мы стремились в статье избежать введения дополнительной


терминологии и дать непосредственные доказательства, не ссылаясь на


специальные результаты о разложении пространств с мерой и булевых


алгебр.


\end{note}





Введенные в этом параграфе определения проиллюстрируем теперь на


простых примерах. В первых трех из них $T=[0,1]$, $\Sigma$ ---


$\sigma$-ал\-геб\-ра измеримых по Лебегу множеств в $T$, $\mu$


--- мера Лебега.





\begin{exam}


$\Sigma_1=\left\{\mathop{\cup}\limits_{n=1}^{m}


\left\{A+\frac{n-1}{m}\right\}\:\Big |\:


A\in\Sigma\cap\left[0,\frac{1}{m}\right]\right\}$ (натуральное


$m>1$ фиксировано).


Очевидно, $\Sigma_1\in{\Omega}(\Sigma)$, более


того, $T_d=T$, $T_c=T_s=\emptyset$ и для


$T_n=\big[\frac{n-1}{m},\frac{n}{m}\big]$ имеет место


$\{T_n\}_{n=1}^m\in\Re(\Sigma_1)$.


\end{exam}





\begin{exam}


$\Sigma_1=\big\{\{t\in T\mid


\cos\frac{\pi}{t}\in V\}\,\big |\,V\in\Sigma


\big\}$. Как и в примере 1, имеем $\Sigma_1\in{\Omega}(\Sigma)$,


$T_d=T$, $T_c=T_s=\emptyset$, и для


$T_n=\big(\frac{1}{n+1},\frac{1}{n}\big]$ справедливо


$\{T_n\}_{n=1}^\infty\in\Re(\Sigma_1)$.


\end{exam}





\begin{exam}


$\Sigma_1$ --- $\sigma$-ал\-геб\-ра, порожденная семейством


$(\Sigma\cap[0,a])\cup\Sigma_0\cup\{T\}$,


где $a\in (0,1)$. В этом случае $\Sigma_1\notin{\Omega}(\Sigma)$,


поскольку $T_c=[0,a]$, $T_s=[a,1]$, $T_d=\emptyset$.


\end{exam}





\begin{exam} $T=[0,1]^2$, $\Sigma$, $\wt{\Sigma}$


--- $\sigma$-ал\-геб\-ры измеримых по Лебегу множеств в $[0,1]^2$


и $[0,1]$ соответственно,


$\Sigma_1=\{A\times[0,1]\mid A\in\wt{\Sigma}\}$.


Тогда $\Sigma_1\notin{\Omega}(\Sigma)$, $T_s=T$,


$T_c=T_d=\emptyset$.


\end{exam}





Очевидно, в примерах 1, 2\; $\sigma$-подалгебры


$\Sigma_1\in\Omega(\Sigma)$ порождаются функциями,


удовлетворяющими условию $\omega(h_1,h_2,\dotsc,h_n)$ работы


(\cite{20}, с.\,97). Читатель без труда сможет найти эти функции


и соответствующие им функции $h_k$.





Через $\wt{\mu}$ обозначим сужение меры $\mu$ на


$\sigma$-ал\-геб\-ру $\Sigma_1$, а через $\wt{\mu}_E$ ($E\in\Sigma$) ---


сужение внешней меры $\wt{\mu}^\ast$ на $\sigma$-ал\-геб\-ру


$\Sigma\cap E$.





\begin{lem} Если $\Sigma_1\in{\Omega}(\Sigma)$,


$R=\{T_n\}_{n\in I}\in\Re(\Sigma_1)$, то при каждом $n\in I$


функция множества $\wt{\mu}_{T_n}$ есть мера, абсолютно


непрерывная относительно меры $\mu_{T_n}:=\mu|_{\Sigma\cap T_n}$.


\end{lem}





\begin{pf} Из равенства (\ref{5}) вытекает счетная


аддитивность функции $\wt{\mu}_{T_n}$, а из леммы~4 ---


свойство $[A\in\Sigma\cap T_n,\,\wt{\mu}_{T_n}A=0]\Rightarrow


[{\mu}_{T_n}A=0]$.


\end{pf}





\begin{note} 1) При $\Sigma_1\notin{\Omega}(\Sigma)$,


$E\in\Sigma\setminus\Sigma_1$ функция множества


$\wt{\mu}_E$ может не быть счетно аддитивной. Это


справедливо, в частности, для $E=[a,1]$ из примера~3 или для


$E=\{(t,s)\in[0,1]^2\mid t\le s\}$ из примера~4.





2) Ecли в условиях леммы 8\; $n\in N$ таково, что $T_n\ne T_c$,


то мера $\wt{\mu}_{T_n}$ не совпадает с $\mu_{T_n}$. Это легко видеть


на примерах~1, 2, в частности, для $T_n$ из примера~1 имеем


$\wt{\mu}_{T_n}A=m\cdot\mu_{T_n}$.


\end{note}





\section{Критерий представимости локального оператора\protect\\ в виде


оператора Немыцкого}





Пусть, как и ранее, $(T,\Sigma,\mu)$ --- пространство с полной


конечной мерой, a


$(X,d)$, $(Y,\rho)$ --- полные сепарабельные метрические пространства.


Вплоть до теоремы~6 будем предполагать,


что $\Sigma_1\in{\Omega}(\Sigma)$,


$R=\{T_n\}_{n\in I}\in\Re(\Sigma_1)$.


Через $\gamma_n$, $r_n$ будем


обозначать операторы, определенные как $\gamma$ в лемме~4 и $r_E$


в теореме~2, где роль $E$ играет $T_n$. Для удобства положим также


$P_n:= P_{T_n,X}$. Через $\wt{\mu}$ обозначим сужение меры $\mu$ на


$\sigma$-ал\-геб\-ру $\Sigma_1$, а через $\mu_n$ --- сужение


внешней меры


$\wt{\mu}^\ast$ на $\sigma$-ал\-геб\-ру $\Sigma\cap T_n$ (т.\,е.


$\mu_n=\wt{\mu}_{T_n}$).





\begin{lem}


$L_0(\Sigma,X)=S_R(L_0(\Sigma_1,X))$.


\end{lem}





\begin{pf} Из включений $T_n\in\Sigma$,


$L_0(\Sigma_1,X)\subset L_0(\Sigma,X)$ следует согласно лемме~2\;


$S_R(L_0(\Sigma_1,X))\subset S_R(L_0(\Sigma,X))=L_0(\Sigma,X)$.





Обратно, зафиксируем произвольное $\f\in L_0(\Sigma,X)$. В силу


теоремы~2 (роль $E$ здесь играют $T_n$, а роль $\Sigma$ --- сначала


$\Sigma$, а затем $\Sigma_1$) и определения $\Omega$-усло\-вия


$$


P_n(L_0(\Sigma,X))=L_0(\Sigma\cap T_n,X)=L_0(\Sigma_1\cap T_n,X)=


P_n(L_0(\Sigma_1,X)), \quad n\in I.


$$


Отсюда следует $L_0(\Sigma,X)\subset S_R(L_0(\Sigma_1,X))$ 


согласно лемме~1.


\end{pf}





\begin{thm}


Предположим, что $(T,\Sigma,\mu)$ ---


пространство с полной конечной мерой, полная $\sigma$-подал\-геб\-ра


$\Sigma_1\subset\Sigma$ удовлетворяет $\Omega$-ус\-ло\-вию,


$(X,d)$, $(Y,\rho)$


--- сепарабельные метрические пространства, и $h:L_0(\Sigma_1,X)\to


L_0(\Sigma,Y)$ --- ЛО.





Тогда существует единственный ЛО $\wt{h}:L_0(\Sigma,X)\to


L_0(\Sigma,Y)$, сужение которого на $L_0(\Sigma_1,X)$ совпадает с $h$.


\end{thm}





Доказательство вытекает непосредственно из теоремы 1 и леммы 9.





Далее ЛО $\wt{h}$, для которого справедливо утверждение теоремы 4,


будем называть


локальным продолжением оператора $h$ на пространство $L_0(\Sigma,X)$.





\begin{thm}


Если в условиях теоремы $4$ ЛО


$h:L_0(\Sigma_1,X)\to L_0(\Sigma,Y)$ непрерывен, то


его локальное продолжение


$\wt{h}:L_0(\Sigma,X)\to L_0(\Sigma,Y)$ есть также непрерывный


оператор.


\end{thm}





\begin{pf} Предположим, что утверждение не имеет


места. Докажем сначала существование таких


$\sigma >0$, $m\in I$ и $\f_n\in L_0(\Sigma,X)$, $n=0,1,\dotsc$, что


$\f_n@>{\mu}>>\f_0$, но


\begin{equation}


\label{8}


\mu E_n >\sigma,


\end{equation}


где


$$


E_n:=\{t\in T_m\mid d((\wt{h}\f_n)(t),(\wt{h}\f_0)(t))>\sigma\}.


$$





Так как оператор $\wt{h}$ не является непрерывным в


точке $\f_0$, то для некоторого $\delta >0$ и некоторой


последовательности $\{\psi_k\}$, сходящейся по мере к $\f_0$, имеем


$\mu D_k>\delta$, где


$$


D_k:=\{t\in


T\mid \rho((\wt{h}\psi_k)(t),(\wt{h}\f_0)(t))>\delta


\}, \quad k=1,2,\dotsc


$$


Из сепарабельности $Y$ вытекает $D_k\in\Sigma$ (см. лемму 2 работы


\cite{78}).


Найдем такое $M\in I$, что


$\mu\mathop{\mathop{\cup}\limits}\limits_{k>M}


T_k<\delta /2$.


Тогда


$\mu\left(\mathop{\cup}\limits_{j=1}^M(D_k\cap T_j)\right)>\delta /2$,


$k\in I$. Отсюда следует существование подпоследовательности


$\{k_n\}$ и индекса $m\in\{1,2,\dotsc,M\}$, для которых


$\mu(D_{k_n}\cap


T_m)>\frac{\delta}{2M}$, $n=1,2,\dotsc$ Тогда неравенство


(\ref{8}) имеет


место для $\f_n=\psi_{k_n}$,


$\sigma=\frac{\delta}{2M}$, $E_n=D_{k_n}\cap T_m$.





Определим функции $\wt{\f}_n\in L_0(\Sigma_1,X)$ равенством


$\wt{\f}_n=r_mP_m\f_n$


($\wt{\f}_n$ действительно принадлежат $L_0(\Sigma_1,X)$,


поскольку $P_m\f_n\in L_0(\Sigma\cap T_m,X)=L_0(\Sigma_1\cap T_m,X)$).


Для произвольного фиксированного $\e >0$ рассмотрим множества


$$


A_n:=\{t\in T_m\mid d(\wt{\f}_n(t),\wt{\f}_0(t))>\e\}, \quad


B_n:=\{t\in T\mid d(\wt{\f}_n(t),\wt{\f}_0(t))>\e\}.


$$


Из сходимости $\f_n@>{\mu}>>\f_0$ и


равенства $P_m\wt{\f_n}=P_m\f_n$ следует


$\mu A_n\to 0$, $n\to\infty$, а из леммы~8 вытекает $\mu_mA_n\to 0$,


$n\to\infty$. Далее, согласно теореме~2\; $B_n=\gamma_m(A_n)$, поэтому


$\mu B_n=\mu_mA_n\to 0$, $n\to\infty$.





Итак, имеем (a) $\wt{\f}_n\in L_0(\Sigma_1,X)$, (b)


$\wt{\f}_n@>{\mu}>>\wt{\f}_0$,


(c) $P_m\f_n=P_m\wt{\f}_n$, $n=1,2,\dotsc$





Из (с), локальности $\wt{h}$ и свойства


$\wt{h}|_{L_0(\Sigma_1,X)}=h$ следует


$$


P_mh\wt{\f}_n=P_m\wt{h}\f_n.


$$


Отсюда и из неравенства (\ref{8}) вытекает, что последовательность


$\{h\wt{\f}_n\}$ не сходится по мере к $\{h\wt{\f}\}$. Это


утверждение и свойства~(a), (b) показывают, что оператор


$h$ не является непрерывным, что противоречит условию


теоремы.


\end{pf}





ЛО $\wt{h}$, для которого справедливо утверждение теоремы 5,


будем называть


локальным непрерывным продолжением оператора $h$ на пространство


$L_0(\Sigma,X)$.





\begin{lem}


Если сепарабельное метрическое пространство $(X,d)$ содержит


хотя бы одно связное множество, отличное от точки, то


существуют такие $x_0\in X$ и


$\alpha>0$, что выполнены свойства


\begin{itemize}


\item[1)] $\{d(x,x_0)\mid x\in X\}\supset[0,\alpha]$,


\item[2)] для любого $u\in(0,\alpha)$ существует такая точка $x\in X$, что в


каждой окрестности точки $x$ найдутся по крайней мере две такие точки


$x_+$ и $x_-$, что $d(x_+,x_0)>u$, $d(x_-,x_0)<u$.


\end{itemize}


\end{lem}





\begin{pf} 1) Согласно условию существует связное множество


$D\subset X$, содержащее по крайней мере две различные точки $x_0$,


$x_\ast$. Пусть $\alpha=d(x_0,x_\ast)$. Как известно (напр.,


\cite{31}, с.\,122), всякая вещественная непрерывная функция на связном


пространстве вместе с двумя различными


значениями принимает все промежуточные значения. Это


справедливо, в частности, для функции $f(\cdot):= d(\cdot,x_0):D\to{\spc


R}$, для которой $f(x_0)=0$, $f(x_\ast)=\alpha$. Таким образом,


$f(D)\supset[0,\alpha]$.





2) Зафиксируем произвольное $u\in(0,\alpha)$ и рассмотрим множества


$$


D_+:= f^{-1}((u,+\infty)), \quad D_0:= f^{-1}(\{u\}), \quad


D_-:= f^{-1}([0,u)).


$$


В силу уже доказанного свойства 1) каждое из этих трех множеств


непусто. Очевидно,


\begin{equation}


\label{9}


\ov{D}_+\subset D_+\cup D_0, \quad


\ov{D}_-\subset D_-\cup D_0, \quad


\ov{D}_+\cup\ov{D}_-=D


\end{equation}


(здесь и далее черта вверху означает замыкание множества в $D$).





Рассмотрим замкнутые множества $Z_+:=\ov{D}_+\cap D_0$,


$Z_-:=\ov{D}_-\cap D_0$. Заметим, что они непусты.


Действительно, если предположить, например, что $Z_+=\emptyset$, то


получим равенство $\ov{D}_+=D_+$, означающее, что $D_+$ ---


от\-кры\-то-зам\-кну\-тое множество в $D$. Это противоречит связности


$D$ (напр., \cite{31}, с.\,118).





Из связности множества $D$ и (\ref{9}) следует


$Z:= Z_+\cap Z_-=\ov{D}_+\cap\ov{D}_-\ne\emptyset$.


Так как


$Z\subset\ov{D}_+\cap\ov{D}_-\cap D_0$, то множество


$Z$ состоит из


точек, предельных как для множества $D_+$, так и для $D_-$. Таким


образом, любая точка $x\in Z$ будет искомой.


\end{pf}





\begin{defn}


Метрическое пространство называется {\it суслинским}


({\it лузинским}), если оно является непрерывным (соответственно непрерывным


и биективным) образом некоторого полного сепарабельного метрического


пространства.


\end{defn}





Накладываемое ниже ограничение: $X$ --- суслинское метрическое


пространство, содержащее связное множество, отличное от точки,


носит весьма общий характер. Оно выполнено, в частности, для замкнутых


подмножеств с непустой внутренностью сепарабельных банаховых пространств.





Далее считаем, что $\Sigma_1$ --- некоторая


полная $\sigma$-по\-дал\-геб\-ра алгебры $\Sigma$, не накладывая априорно


ограничения $\Sigma_1\in{\Omega}(\Sigma)$.





Приведем для удобства ряд известных определений


(см., напр., \cite{28}--\cite{30}).





\begin{defn}


Пространство с полной конечной мерой $(T,\Sigma, \mu)$ называется


{\it про\-ст\-ран\-ст\-вом Ле\-бе\-га}, если существует


счетная подалгебра ${\cal A}\subset\Sigma$ такая, что выполнены условия


\begin{itemize}


\item[a)] $\Sigma$ есть пополнение по мере $\mu$\; $\sigma$-ал\-геб\-ры,


порожденной $\cal A$,


\item[b)] $\cal A$ разделяет точки $T$, т.\,е.


$$


(\forall t,s\in T,\; t\ne s)\;\ (\exists A\in{\cal A}):\;\ 


(t\in A,\; s\in T\setminus A),


$$


\item[c)] $\cal A$ --- компактный класс множеств, т.\,е. каждая убывающая


последовательность множеств из $\cal A$ имеет непустое пересечение.


\end{itemize}


\end{defn}





\begin{defn}


Измеримые пространства $(T_i,{\cal T}_i)$,


$i=1,2$, называются {\it борелевски изоморфными}, если существует такое


биективное отображение $f$ из $T_1$ на $T_2$, что


$f^{-1}({\cal T}_2)\subset{\cal T}_1$ и


$(f^{-1})^{-1}({\cal T}_1)\subset{\cal T}_2$. При этом отображение $f$


называется {\it борелевским изоморфизмом} $T_1$ на $T_2$.


\end{defn}





\begin{defn}


Пространства с полной конечной мерой $(T_i,{\cal T}_i, \mu_i)$, $i=1,2$,


называются {\it изоморфными $\moddd$}, если существует отображение $f:T_1\to


T_2$ и множества $N_i\subset {\cal T}_i$ с мерой $\mu_i


N_i=0$, $i=1,2$, такие, что выполняются условия


\begin{itemize}


\item[a)] $f|_{T_1\setminus N_1}$ есть борелевский изоморфизм измеримого


пространства $(T_1\setminus N_1,{\cal T}_1\cap(T_1\setminus N_1))$ на


$(T_2\setminus N_2,{\cal T}_2\cap(T_2\setminus N_2))$,


\item[b)] $\mu_1f^{-1}=\mu_2$.


\end{itemize}


При этом отображение $f$ называется {\it изоморфизмом} из пространства


с мерой $(T_1,{\cal T}_1, \mu_1)$ на $(T_2,{\cal T}_2, \mu_2)$.


\end{defn}





Хорошо известна


\medskip





{\bf Теорема Рохлина} (\cite{28}).


{\it $1)$ Пространство, изоморфное пополнению


стандартного пространства $($т.\,е. полного сепарабельного


метрического пространства с борелевской мерой$)$, есть пространство


Лебега.





$2)$ Пространство Лебега $(T,\Sigma,\mu)$ с неатомической мерой


изоморфно $\moddd$ пространству $([0,1],{\cal L},m_T)$, где $\cal L$


--- лебеговская $\sigma$-ал\-геб\-ра на $[0,1]$, а $m_T$


--- мера, равная мере Лебега, умноженной на число $\mu T$ $($так что


$m_T[0,1]=\mu T)$.


}





\begin{defn}


Для функции $f:T\times X\to Y$ оператор ${\cal


N}_f:F(X)\to F(Y)$, определяемый равенством


$$


({\cal N}_f\f)(t)=f(t,\f(t)) \quad \text{при почти всех} \quad


t\in T,


$$


называется {\it оператором Немыцкого (суперпозиции)}, порожденным функцией


$f$.


\end{defn}





\begin{defn}


Функция $f:T\times X\to Y$ называется


{\it функцией Шрагина} (\cite{7}, \cite{8}),


если существует такое $T_0\in\Sigma_0$,


что сужение функции $f$ на множество $(T\setminus T_0)\times X$\;


$\Sigma\otimes{\cal B}(X)$-из\-ме\-ри\-мо


($\Sigma\otimes{\cal B}(X)$ --- $\sigma$-ал\-геб\-ра, порожденная


всевозможными прямоугольниками $A\times B$, где $A\in\Sigma$,


$B\in{\cal B}(X)$).


\end{defn}





Отметим, что в \cite{32} и \cite{55} такая функция названа


стандартной.





\begin{defn}


Говорят, что функция $f:T\times X\to Y$


удовлетворяет {\it условиям Каратеодори}, если для любого $x\in X$ функция


$f(\cdot, x)$\; $\Sigma$-из\-ме\-ри\-ма и при почти всех $t\in T$


функция $f(t,\cdot)$ непрерывна.


\end{defn}





\begin{thm}


Пусть $(T,\Sigma,\mu)$ --- пространство Лебега,


$X$ --- суслинское метрическое пространство, содержащее


связное множество, отличное от точки,


и $Y$ состоит более чем из одной точки.





Тогда, если $\sigma$-подалгебра $\Sigma_1$ не удовлетворяет


$\Omega$-ус\-ло\-вию, то существует непрерывный ЛО


$h:L_0(\Sigma_1,X)\to L_0(\Sigma,Y)$, не продолжаемый локально и


непрерывно на $L_0(\Sigma,X)$ $($т.\,е. для любого непрерывного ЛО


$h^\ast:L_0(\Sigma,X)\to L_0(\Sigma,Y)$ имеет место


$h^\ast|_{L_0(\Sigma_1,X)}\ne h)$.


\end{thm}





{\bf Доказательство.} Согласно следствию 1\; $\Sigma_s\ne\emptyset$,


а согласно утверждению~3) теоремы~3 множество $T_s$ (положительной меры)


обладает свойствами





A) $[A\subset T_s,\,\mu A>0]\,\Rightarrow\,


[\Sigma\cap A\ne\Sigma_1\cap A];$





B) $T_s$ не содержит атомов меры $\mu$.





Согласно теореме Рохлина существует изоморфизм $\beta$


пространства $(T,\Sigma\cap T_s,\mu|_{\Sigma\cap T_s})$ на


$([0,1],{\cal L},m_{T_s})$





Пусть $y_1$, $y_2$ --- два различных элемента метрического


пространства $Y$. Выберем точку $x_0\in X$ и


$\alpha\in(0,1)$ такие, что справедливо утверждение леммы~10.


Определим функцию $f:T\times X\to Y$ равенством


\begin{equation}


\label{10}


f(t,x)=\left\{ \begin{array}{rlcl}


y_1, & \text{если } t\in T\setminus T_s & \text{ или } &


d(x,x_0)\ge\beta(t);\\


y_2, & \text{если } t\in T_s & \text{ и } & d(x,x_0)<\beta(t),


\end{array} \right.


\end{equation}


и рассмотрим оператор Немыцкого ${\cal N}_f$, порожденный функцией $f$.





Легко видеть, учитывая (\cite{78}, лемма 2), что


функция $f$\; $\Sigma\otimes{\cal B}(X)$-из\-ме\-ри\-ма,


тем более является функцией Шрагина. Поэтому,


как доказано в \cite{32}, \cite{55},


${\cal N}_f:L_0(\Sigma,X)\to L_0(\Sigma,Y)$. Далее докажем


последовательно свойства {\bf 1)}, {\bf 2)},


из которых вытекает утверждение теоремы.





{\bf 1)} ${\cal N}_f:L_0(\Sigma_1,X)\to L_0(\Sigma,Y)$ и непрерывен.





{\bf 2)} Не существует непрерывного ЛО


$h^\ast:L_0(\Sigma,X)\to L_0(\Sigma,Y)$, сoвпадающего с ${\cal N}_f$ на


подпространстве $L_0(\Sigma_1,X)$.





Несложное рассуждение ``от про\-тив\-но\-го'' с использованием


теоремы Егорова показывает, что для доказательства {\bf 1)} достаточно


для произвольной последовательности


$\{\f_n\}_{n=1}^\infty\subset L_0(\Sigma_1,X)$,


сходящейся п.\,в. к $\f_0$, доказать, что ${\cal N}_f\f_n\to


{\cal N}_f\f_0$ п.\,в. на $T$. Зафиксируем такую


последовательность и, положив


$$


A:=\{t\in T_s\mid d(\f_0(t),x_0)=\beta(t)\}\cup


\{t\in T\mid \f_n(t)\not\to\f_0(t)\},


$$


покажем, что


\begin{equation}


\label{11}


\mu A=0.


\end{equation}





Предположим противное. Из теоремы 3 следует существование множества


$\wt{E}$ ненулевой меры, принадлежащего семейству $(\Sigma\cap


A)\setminus (\Sigma_1\cap A)$. Из определения функции $\beta$ вытекает


включение $\beta(\wt{E})\in{\cal L}$. Следовательно, существует


множество $E$, равное $\wt{E}\;\moddd$, для которого


$V:=\beta(D)\in{\cal B}([0,1])$.





Из $\Sigma_1$-из\-ме\-ри\-мос\-ти функции $\f_0$ следует $F:=


[d(\f_0(\cdot),x_0)]^{-1}(V)\in\Sigma_1$. Согласно определению


множеств $A$, $E$, $V$, $F$


$$


\Sigma_1\cap A\ni F\cap A=\beta^{-1}(V)\cap A=E.


$$


Это противоречит свойству $E\notin\Sigma_1\cap A$. Таким образом,


равенство (\ref{11}) доказано.





Далее, $(\forall t\in T_s\setminus A)$\; $d(\f_0(t),x_0)\ne\beta(t)$,


поэтому в силу сходимости $\f_n(t)\to\f_0(t)$ имеем при всех


$n$, б\'{о}льших некоторого $N(t)$,


$$


\text{sign}\,(d(\f_n(t),x_0)-\beta(t))=


\text{sign}\,(d(\f_0(t),x_0)-\beta(t)).


$$


Таким образом, последовательность $\{f(t,\f_n(t))\}_{n=N(t)}^\infty$


стационарна. Мы показали, что\linebreak


$({\cal N}_f\f_n)(t)\to({\cal N}_f\f_0)(t)$.


Отсюда и из (\ref{11}), (\ref{12})


следует сходимость ${\cal N}_f\f_n\to{\cal N}_f\f_0$


п.\,в. на $T$. Свойство {\bf 1)} доказано.





Предположим, что свойство {\bf 2)} не имеет места. Согласно теореме о


представлении \cite{4}--\cite{6} существует функция $g:T\times X\to


Y_c$ ($(Y_c,\rho_c)$ --- пополнение метрического пространства


$(Y,\rho)$), удовлетворяющая условиям Каратеодори и порождающая


оператор $h^\ast$ (таким образом, $h^\ast={\cal N}_g$).





Положим


$$


\wt{T}=\{t\in T_s\cap


\beta^{-1}((0,1))\mid f(t,\cdot)\; \text{ непрерывна на }\;X\}.


$$


Очевидно, $\mu\wt{T}>0$. Зафиксируем произвольное счетное всюду


плотное в $X$ множество $\{x_n\}_{n=1}^\infty$ и рассмотрим множества


\begin{equation}


\label{12}


E_n:=\{t\in\wt{T}\mid f(t,x_n)\ne g(t,x_n)\}


\end{equation}


(точки $Y$ отождествляем с точками $Y_c$ при естественном вложении).


Очевидно, $(\forall n)$\; $f(\cdot,x_n)$, $g(\cdot,x_n)\in L_0(\Sigma,


Y_c)$, поэтому $E_n\in\Sigma$. Покажем, что


\begin{equation}


\label{13}


\bigcup_{n=1}^\infty E_n=\wt{T}.


\end{equation}





Зафиксируем произвольное $t\in\wt{T}$ и найдем согласно


утверждению 2) леммы 10 такое $x\in X$, что


\begin{equation}


\label{14}


B_r\cap X_1\ne\emptyset, \quad


B_r\cap X_2\ne\emptyset \quad (\forall r>0),


\end{equation}


где


$$


B_r:=\{\xi\in X\mid d(\xi,x)<r\}, \quad


X_1:=\{\xi\in X\mid d(\xi,x_0)>\beta(t)\}, \quad


X_2:=\{\xi\in X\mid d(\xi,x_0)<\beta(t)\}.


$$





В силу непрерывности функции $g(t,\cdot)$ существует такое $\delta>0$,


что


\begin{equation}


\label{15}


\sup_{\xi\in B_\delta}\rho_c(g(t,\xi),g(t,x))<


\rho(y_1,y_2)/2.


\end{equation}





Ввиду открытости множеств $B_\delta\cap X_i$, $i=1,2$, и свойства (\ref{14})


существуют такие натуральные числа $n_i$, что $x_{n_i}\in B_\delta\cap


X_i$, $i=1,2$. Согласно (\ref{10}) $f(t,x_{n_i})=y_i$, $i=1,2$. Отсюда и из


(15) следует, что либо для $i=1$, либо для $i=2$ имеет место


$f(t,x_{n_i})\ne g(t,x_{n_i})$. Таким образом, равенство (\ref{13}) доказано.





Из (\ref{13}) и свойства $\mu\wt{T}>0$ следует, что для некоторого


натурального $m$\; $\mu E_m>0$. Тогда для функции $\f\in


L_0(\Sigma_1,X)$ вида $\f(t)\equiv x_m$ получаем в силу (\ref{12})


неравенство ${\cal N}_f\f\ne{\cal N}_g\f$. Таким образом, сужения


операторов ${\cal N}_f$ и ${\cal N}_g$ на $L_0(\Sigma_1,X)$ не совпадают.


Полученное противоречие доказывает утверждение {\bf 2)}.


%$\quad\Box$





Из теорем 5, 6 и теоремы о представлении [4]--[6]


вытекает основной результат статьи.





\begin{thmmain}


Пусть $(T,\Sigma,\mu)$ --- пространство Лебега, $\Sigma_1$ ---


полная $\sigma$-подал\-гебра $\sigma$-ал\-геб\-ры $\Sigma$,


$X$ --- суслинское метрическое пространство, содержащее


связное множество, отличное от точки, и $Y$ ---


лузинское метрическое пространство, состоящее более чем из одной


точки.





Для того чтобы любой непрерывный ЛО 


$h:L_0(\Sigma_1,X)\to L_0(\Sigma,Y)$ был представим в виде оператора


Немыцкого с порождающей функцией, удовлетворяющей условиям Каратеодори,


необходимо и достаточно, чтобы $\sigma$-подал\-гебра $\Sigma_1$


удовлетворяла $\Omega$-ус\-ло\-вию $($см. $(5))$.


\end{thmmain}





\begin{note}


Утверждение теоремы 7


является, по-видимому, новым даже для простейшего случая


$T=[0,1]$ (с мерой Лебега), $X=Y={\spc R}$.


\end{note}





\begin{note}


Вопрос о том, при каких ограничениях на полную


$\sigma$-подал\-геб\-ру $\Sigma_1$ любой непрерывный ЛО


$h:L_0(\Sigma_1,X)\to L_0(\Sigma,Y)$ представим в виде оператора


Немыцкого с порождающей функцией Шрагина, в настоящее время открыт.


Здесь интерес представляет исследование ``узких'' $\sigma$-подал\-гебр,


не удовлетворяющих ${\Omega}$-ус\-ло\-вию (для $\Sigma_1\in


{\Omega}(\Sigma)$ положительный ответ очевидным образом вытекает из теоремы


7 и того факта, что любая функция Каратеодори есть функция Шрагина


\cite{32}).


\end{note}





Теоремы 5--7 позволяют в формулируемом ниже следствии


полностью описать связь между


$\Omega$-ус\-ло\-ви\-ем на $\sigma$-подал\-гебру $\Sigma_1$,


свойством непрерывной продолжаемости ЛО


и фактом представимости оператора


в виде оператора Немыцкого с порождающей функцией Каратеодори.





Через ${\cal N}(\Sigma)$ обозначим класс таких полных $\sigma$-подал\-гебр


алгебры $\Sigma$, что для любого непрерывного ЛО


$h:L_0(\Sigma_1,X)\to L_0(\Sigma,Y)$ существует локальное непрерывное


продолжение $h$ на $L_0(\Sigma, X)$.





Через ${\cal K}(\Sigma)$ обозначим класс таких полных $\sigma$-подал\-гебр


алгебры $\Sigma$, что любой непрерывный ЛО


$h:L_0(\Sigma_1,X)\to L_0(\Sigma,Y)$ есть оператор Немыцкого с


порождающей функцией, удовлетворяющей условиям Каратеодори.





\begin{cor}


1. ${\cal K}(\Sigma)\subset{\cal N}(\Sigma)\subset{\Omega}(\Sigma)$.





2. Если $(T,\Sigma,\mu)$ --- пространство Лебега, а метрические


пространства $X$, $Y$ удовлетворяют условиям теоремы 7, то


${\cal K}(\Sigma)={\cal N}(\Sigma)={\Omega}(\Sigma)$.


\end{cor}





\begin{thebibliography}{99}





\bibitem{1} Шрагин И.В. {\it Абстрактные операторы Немыцкого --- локально


определенные операторы} // ДАН СССР. -- 1976. -- Т.\,227. -- \No\,1. --


С.\,47--49.





\bibitem{2} Kart\'{a}k K. {\it A generalization of the Carath\'{e}odory


of differential equations} // Czechosl. Math. J. -- 1967. -- Т.\,17. --


\No\,4. -- P.\,482--514.





\bibitem{3} Vrko\u{c} I. {\it The representation of Carath\'{e}odory


operators} // Czechosl. Math. J. -- 1969. -- Т.\,19. -- \No\,1. --


P.\,99--109.





\bibitem{4} Поносов А.В. {\it Операторы Каратеодори и операторы Немыцкого}.


-- Пермск. ун-т. -- Пермь, 1984. -- 18~с. -- Деп. в ВИНИТИ 18.07.84,


\No\,5141--84.





\bibitem{5} Поносов А.В. {\it К теории локально определенных


операторов} // Функц.-дифференц. уравнения. -- Пермь,


1985. -- С.\,72--82.





\bibitem{6} Поносов А.В. {\it О гипотезе Немыцкого} // ДАН СССР. --


1986. -- Т.\,289. -- \No\,6. -- С.\,1308--1311.





\bibitem{7} Appell J. {\it The superposition operator in function spaces


--- A survey} // Expos. Math. -- 1988. -- V.\,6. -- \No\,3. -- P.\,209--270.





\bibitem{8} Appell J., Zabrejko P.P. {\it Nonlinear superposition


operators}. -- Cambridge: Cambr. Univ. Press, 1990. -- 312 p.





\bibitem{9} Красносельский М.А. {\it Топологические методы в теории


нелинейных интегральных уравнений}. -- М.: Гос\-тех\-из\-дат, 1956. -- 392~с.





\bibitem{21} Вайнберг М.М. {\it Вариационные методы исследования


нелинейных операторов}. -- М.: Гостехиздат, 1956. -- 344~с.





\bibitem{10} Красносельский М.А., Забрейко П.П., Пустыльник Е.И.,


Соболевский П.Е.


{\it Интегральные операторы в пространствах суммируемых функций}. --


М.: Наука, 1966. -- 499~с.





\bibitem{32} Шрагин И.В. {\it Условия измеримости суперпозиций} // ДАН СССР.


-- 1971. -- Т.\,197. -- \No\,2. -- С.\,295--298.





\bibitem{55} Шрагин И.В. {\it Суперпозиционная измеримость} // Изв. вузов.


Математика. -- 1975. -- \No\,1. -- С.\,82--92.





\bibitem{26} Красносельский М.А., Покровский А.В. {\it О разрывном операторе


суперпозиции} // УМН. -- 1977. -- Т.\,32. -- Вып.\,1. --


С.\,169--170.





\bibitem{78} Шрагин И.В. {\it Необходимость условий Каратеодори для


непрерывности оператора Немыцкого} //


Функц.-дифференц.


уравнения и краев. задачи матем. физики. -- Пермь, 1978. -- С.\,128--134.





\bibitem{11} Шрагин И.В. {\it Локально определенные операторы и гипотеза


Немыцкого} //


Функц.-дифференц. уравнения. -- Пермь, 1991. --


С.\,95--101.





\bibitem{22} Красносельский М.А., Покровский А.В.


{\it Системы с гистерезисом}. -- М.: Наука, 1993. -- 272~с.





\bibitem{12} Шрагин И.В. {\it Классы измеримых век\-тор-функ\-ций и оператор


Немыцкого}. I. // Изв. вузов. Математика. -- 1994. -- \No\,4. -- С.\,48--58.





\bibitem{99} Шрагин И.В. {\it Классы измеримых век\-тор-функ\-ций и оператор


Немыцкого}. II. // Изв. вузов. Математика. -- 1994. --


\No\,5. -- С.\,70--79.





\bibitem{23} Shragin I.V. {\it On representation of a locally


defined operator in the form of the Nemytskii operator} //


Funct. Different. Equat. Israel Seminar. -- 1996. --


\No\,3--4. -- P.\,447--452.





\bibitem{13} Шрагин И.В., Непомнящих Ю.В. {\it $D$-ус\-ло\-вия


Каратеодори и их


связь с $D$-не\-пре\-рыв\-ностью оператора Немыцкого} //


Изв. вузов. Математика. -- 1997. -- \No\,6. -- С.\,70--82.





\bibitem{14} Поносов А.В. {\it Теорема о неподвижной точке для


локально определенных операторов} //


Краев. задачи. -- Пермь, 1985.


-- С.\,125--128.





\bibitem{15} Поносов А.В. {\it Метод неподвижной точки в теории


стохастических дифференциальных уравнений} // ДАН СССР. -- 1988. --


Т.\,299. -- \No\,3. -- С.\,562--565.





\bibitem{16} Поносов А.В. {\it К теории приводимых


функ\-ци\-о\-наль\-но-диф\-фе\-рен\-ци\-аль\-ных уравнений Ито} //


Дифференц. уравнения. -- 1989. -- Т.\,25. -- \No\,11. -- С.\,1915--1925.





\bibitem{20} Драхлин М.Е.


{\it Об одном линейном функциональном уравнении} //


Функц.-дифференц. уравнения. -- Пермь,


1985. -- С.\,91--111.





\bibitem{33} Азбелев Н.В., Максимов В.П., Рахматуллина Л.Ф.


{\it Введение в теорию функ\-ци\-о\-наль\-но--диф\-фе\-рен\-ци\-аль\-ных


уравнений.} -- М.: Наука, 1991. -- 280~с.





\bibitem{34} Azbelev N.V. {\it The ideas and methods of the Perm


Seminar on boundary value problems} // Boundary Value Problems


for Functional Different. Equat. -- Singapore: World Scientific


Publishing, 1995. -- P.\.13--22.





\bibitem{24} Харазишвили А.Б. {\it Приложения теории множеств.} -- Тбилиси:


Изд-во Тбилис. ун-та, 1989. -- 143~с.





\bibitem{25} Владимиров Д.А. {\it Булевы алгебры}. --


М.: Наука, 1969. -- 318~с.





\bibitem{27} Халмош П.Р. {\it Теория меры.} -- М.: Ин. лит., 1953. --


292~с.





\bibitem{31} Александров П.С. {\it Введение в теорию множеств и


общую топологию}. Учеб. пособие. -- М.: Наука, 1977. -- 368~с.





\bibitem{28} Рохлин В.А. {\it Об основных понятиях теории меры} //


Матем. сб. -- 1949. -- Т.\,25. -- \No\,1. -- С.\,107--150.





\bibitem{29} Самородницкий А.А. {\it Теория меры.}


-- Л.: Изд-во ЛГУ, 1990. -- 267~с.





\bibitem{30} Партасарати К.


{\it Введение в теорию вероятностей и теорию меры.}


-- М.: Мир, 1983. -- 343~с.





\end{thebibliography}





\vskip 2cm





\post{\it Пермский государственный университет,}{\it Поступила}


\post{Университет NLH, Норвегия}{18.03.1998}





\end{document}





